% !TEX root = ../../ctfp-print.tex

\lettrine[lhang=0.17]{范}{畴论中的}大多数构造, 都是对数学中其他更具体领域的结果的推广. 积、余积、
幺半群、指数这些东西, 早在范畴论诞生之前就已为人所知. 它们在数学的不同分支里或许有着不同的名字.
集合论中的笛卡尔积、序理论中的交 (meet)、逻辑中的合取 --- 它们都是范畴论中积这个抽象想法的具体例子.

米田引理在这方面格外出众: 它是一个关于范畴总体的、覆盖面极广的论断, 在其他数学分支里几乎没有先例.
有人说, 与它最接近的类比是群论中的 Cayley 定理 (每个群都同构于某个集合的置换群).

米田引理的舞台是任意一个范畴 $\cat{C}$, 加上一个从 $\cat{C}$ 到 $\Set$ 的函子 $F$.
我们在上一节看到, 某些取值于 $\Set$ 的函子是可表的, 也就是说它们同构于某个 hom-函子.
米田引理告诉我们: 所有取值于 $\Set$ 的函子都能通过自然变换从 hom-函子得到, 而且它还明确地
列举出了所有这样的变换.

我讲自然变换的时候提过, 自然性条件可能相当苛刻. 当你在某个物件处定义好自然变换的一个分量,
自然性可能强到足以把这个分量 ``搬运'' 到另一个与它有态射相连的物件上. 源范畴与目标范畴的物件之间
的箭头越多, 你在搬运自然变换的分量时受到的约束就越多. 而 $\Set$ 恰好是一个箭头极其丰富的范畴.

米田引理告诉我们: hom-函子与任何其他函子 $F$ 之间的自然变换, 只需指定它唯一那个分量在单个点上的
取值, 就完全确定了! 自然变换的其余部分全都由自然性条件推出.

那么先来回顾一下米田引理涉及的两个函子之间的自然性条件. 第一个函子是 hom-函子. 它把 $\cat{C}$
中的任意物件 $x$ 映到态射的集合 $\cat{C}(a, x)$ --- 其中 $a$ 是 $\cat{C}$ 中一个固定的物件.
我们也看到过, 它把 $x \to y$ 的任意态射 $f$ 映到 $\cat{C}(a, f)$.

第二个函子是任意的取值于 $\Set$ 的函子 $F$.

我们把这两个函子之间的自然变换叫作 $\alpha$. 因为我们在 $\Set$ 中操作, 自然变换的分量 --- 比如
$\alpha_x$ 或 $\alpha_y$ --- 不过是集合之间的普通函数:
\begin{gather*}
  \alpha_x \Colon \cat{C}(a, x) \to F x \\
  \alpha_y \Colon \cat{C}(a, y) \to F y
\end{gather*}

\begin{figure}[H]
  \centering
  \includegraphics[width=0.4\textwidth]{images/yoneda1.png}
\end{figure}

\noindent
而这些分量既然就是函数, 我们就可以查看它们在特定点上的取值. 但集合 $\cat{C}(a, x)$ 中的一个
点是什么? 这里有一个关键观察: 集合 $\cat{C}(a, x)$ 中的每个点, 同时也是一个从 $a$ 到 $x$
的态射 $h$.

于是 $\alpha$ 的自然性方形:
\[\alpha_y \circ \cat{C}(a, f) = F f \circ \alpha_x\]
逐点地、当它作用在 $h$ 上时, 变成:
\[\alpha_y (\cat{C}(a, f) h) = (F f) (\alpha_x h)\]
你或许还记得上一节中, hom-函子 $\cat{C}(a,-)$ 对态射 $f$ 的作用被定义为前组合:
\[\cat{C}(a, f) h = f \circ h\]
这就导出:
\[\alpha_y (f \circ h) = (F f) (\alpha_x h)\]
这个条件到底有多强, 只要把它特化到 $x = a$ 的情形就能看出来.

\begin{figure}[H]
  \centering
  \includegraphics[width=0.4\textwidth]{images/yoneda2.png}
\end{figure}

\noindent
那种情形下 $h$ 成了一个从 $a$ 到 $a$ 的态射. 我们知道至少存在这样一个态射, 就是
$h = \id_a$. 把它代进去:
\[\alpha_y f = (F f) (\alpha_a \id_a)\]
留意刚刚发生了什么: 等式左边是 $\alpha_y$ 作用在 $\cat{C}(a, y)$ 的任意元素 $f$ 上. 而它完全由
$\alpha_a$ 在 $\id_a$ 处的唯一一个取值决定. 我们可以任选这样一个值, 它就会生成一个自然变换.
由于 $\alpha_a$ 的取值落在集合 $F a$ 里, $F a$ 中的任意一个点都能定义出某个 $\alpha$.

反过来, 给定任意一个从 $\cat{C}(a, -)$ 到 $F$ 的自然变换 $\alpha$, 你可以在 $\id_a$ 处求值,
得到 $F a$ 中的一个点.

我们刚刚证明了米田引理:

\begin{quote}
  从 $\cat{C}(a, -)$ 到 $F$ 的自然变换, 与 $F a$ 的元素之间存在一一对应.
\end{quote}
换句话说,
\[\cat{Nat}(\cat{C}(a, -), F) \cong F a\]
或者, 如果我们用记号 $[\cat{C}, \Set]$ 表示 $\cat{C}$ 与 $\Set$ 之间的函子范畴, 那么自然变换的
集合不过是这个范畴里的一个 hom-集, 于是我们可以写成:
\[[\cat{C}, \Set](\cat{C}(a, -), F) \cong F a\]
稍后我会解释, 这个对应关系实际上是一个自然同构.

现在我们来找找这个结果的直觉. 最不可思议的是, 整个自然变换只从一个成核点结晶而出: 就是我们在
$\id_a$ 处赋予它的那个取值. 它从那个点出发, 沿着自然性条件蔓延开去, 漫过 $\cat{C}$ 在 $\Set$
中的像. 所以先来分析一下, 在 $\cat{C}(a, -)$ 之下 $\cat{C}$ 的像是什么.

先从 $a$ 自身的像说起. 在 hom-函子 $\cat{C}(a, -)$ 之下, $a$ 被映到集合 $\cat{C}(a, a)$.
而在函子 $F$ 之下, 它被映到集合 $F a$. 自然变换的分量 $\alpha_a$ 是从 $\cat{C}(a, a)$
到 $F a$ 的某个函数. 我们只关注集合 $\cat{C}(a, a)$ 中的一个点, 对应于态射 $\id_a$ 的那个点.
为了强调它不过是集合里一个点, 就叫它 $p$ 吧. 分量 $\alpha_a$ 应该把 $p$ 映到 $F a$ 中的某个点
$q$. 我要向你说明: 无论怎么选 $q$, 都通向唯一的自然变换.

\begin{figure}[H]
  \centering
  \includegraphics[width=0.3\textwidth]{images/yoneda3.png}
\end{figure}

\noindent
第一个论断是: 一个点 $q$ 的选择, 唯一地决定了函数 $\alpha_a$ 的其余部分. 的确, 让我们任取
$\cat{C}(a, a)$ 中的另一个点 $p'$, 它对应于某个从 $a$ 到 $a$ 的态射 $g$. 米田引理的魔法就在此处
上演: $g$ 既可以被看作集合 $\cat{C}(a, a)$ 中的一个点 $p'$, 与此同时, 它又挑选出了集合之间的
两个\emph{函数}. 的确, 在 hom-函子之下, 态射 $g$ 被映到一个函数 $\cat{C}(a, g)$; 而在 $F$ 之下,
它被映到 $F g$.

\begin{figure}[H]
  \centering
  \includegraphics[width=0.4\textwidth]{images/yoneda4.png}
\end{figure}

\noindent
现在考虑 $\cat{C}(a, g)$ 作用在我们最初那个 $p$ 上的情形 --- 你大概还记得, 它对应于 $\id_a$.
这个作用定义为前组合 $g \circ \id_a$, 它等于 $g$, 而 $g$ 对应于我们的点 $p'$. 于是态射
$g$ 被映到一个函数上, 这个函数作用在 $p$ 上产生 $p'$, 而 $p'$ 就是 $g$. 我们绕了一整圈回到
原点了!

再考虑 $F g$ 作用在 $q$ 上. 结果是某个 $q'$, $F a$ 中的一个点. 要让自然性方形闭合, $p'$
必须在 $\alpha_a$ 之下被映到 $q'$. 我们挑了一个任意的 $p'$ (任意的 $g$), 并推出了它在
$\alpha_a$ 之下的映射. 这样一来, 函数 $\alpha_a$ 就完全确定了.

第二个论断是: 对 $\cat{C}$ 中任何与 $a$ 有连接的物件 $x$, $\alpha_x$ 都是唯一确定的. 推理
方式类似, 只不过现在我们多了两个集合, $\cat{C}(a, x)$ 和 $F x$, 而从 $a$ 到 $x$ 的态射 $g$
在 hom-函子之下被映到:
\[\cat{C}(a, g) \Colon \cat{C}(a, a) \to \cat{C}(a, x)\]
在 $F$ 之下被映到:
\[F g \Colon F a \to F x\]
同样, $\cat{C}(a, g)$ 作用在我们的 $p$ 上由前组合 $g \circ \id_a$ 给出, 它对应于
$\cat{C}(a, x)$ 中的一个点 $p'$. 自然性把 $\alpha_x$ 作用在 $p'$ 上的值确定为:
\[q' = (F g) q\]
由于 $p'$ 是任取的, 整个函数 $\alpha_x$ 就这样确定了.

\begin{figure}[H]
  \centering
  \includegraphics[width=0.4\textwidth]{images/yoneda5.png}
\end{figure}

\noindent
如果 $\cat{C}$ 中有与 $a$ 毫无连接的物件呢? 它们在 $\cat{C}(a, -)$ 之下全被映到同一个集合
--- 空集. 回想一下, 空集是集合范畴里的初始物件. 这意味着从这个集合到任何别的集合存在唯一的
函数. 我们把这个函数叫作 \code{absurd}. 所以在这里, 我们同样没有选择的余地: 自然变换的那个分量
只能是 \code{absurd}.

理解米田引理的另一种方式, 是认识到: 取值于 $\Set$ 的函子之间的自然变换不过是函数的族, 而函数
一般是有损的. 一个函数可能折叠信息, 也可能只覆盖它陪域的一部分. 唯一无损的函数是可逆的那些 ---
即同构. 由此可见, 最能保持结构的取值于 $\Set$ 的函子就是可表的那些. 它们要么是 hom-函子, 要么
是与某个 hom-函子自然同构的函子. 任何其他函子 $F$ 都是从 hom-函子经过一次有损变换得到的.
这样的变换不仅可能丢失信息, 还可能只覆盖函子 $F$ 在 $\Set$ 中的像的一小部分.

\section{Haskell 中的米田引理}

我们已经在 Haskell 里遇见过 hom-函子, 它披着 reader 函子的外衣:

\src{snippet01}
reader 通过前组合来映射态射 (在这里, 就是函数):

\src{snippet02}
米田引理告诉我们, reader 函子可以自然地映射到任何别的函子.

自然变换就是一个多态函数. 所以给定一个函子 \code{F}, 我们就有一个从 reader 函子到它的映射:

\src{snippet03}
照例, \code{forall} 可以省略, 但我喜欢把它明确写出来, 以强调自然变换的参数化多态性.

米田引理告诉我们, 这些自然变换与 \code{F a} 的元素之间存在一一对应:

\begin{snipv}
forall x . (a -> x) -> F x \ensuremath{\cong} F a
\end{snipv}
这个等式的右边, 正是我们通常认为的数据结构. 还记得把函子解释成广义容器吗? \code{F a} 是一个装着
\code{a} 的容器. 而左边是一个以函数为参数的多态函数. 米田引理告诉我们, 这两种表示是等价的 ---
它们包含同样的信息.

换一种说法是: 给我一个类型为:

\src{snippet04}
的多态函数, 我就给你一个装 \code{a} 的容器. 诀窍正是我们在米田引理的证明里用过的那个: 我们用
\code{id} 调用这个函数, 得到一个 \code{F a} 的元素:

\src{snippet05}
反过来也成立: 给定一个类型为 \code{F a} 的值:

\src{snippet06}
就可以定义一个类型正确的多态函数:

\src{snippet07}
你可以在这两种表示之间自如往返.

拥有多种表示的好处在于: 一种可能比另一种更容易组合, 或者在某些应用里比另一种更高效.

对这条原则最简单的例证, 就是编译器构造中常用的那种代码变换: 续体传递风格 (continuation passing style), 即 \acronym{CPS}.
它是把米田引理应用于恒等函子的最简单的情形. 把 \code{F} 换成恒等就得到:

\begin{snipv}
forall r . (a -> r) -> r \ensuremath{\cong} a
\end{snipv}
这个公式的解释是: 任何类型 \code{a} 都可以替换成一个接受 \code{a} 的 ``处理器'' 的函数.
处理器是一个接受 \code{a} 并执行计算的剩余部分的函数 --- 即续体. (类型 \code{r} 通常封装了
某种状态码.)

这种编程风格在 UI、异步系统和并发编程中非常常见. \acronym{CPS} 的缺点在于它涉及控制反转.
代码被拆分在生产者与消费者 (处理器) 之间, 不易组合. 任何做过一些不那么简单的 web 编程的人都熟悉
那种噩梦: 有状态的处理器互相作用, 缠成一团面条代码. 我们稍后会看到, 恰当地使用函子和单子,
可以恢复 \acronym{CPS} 的某些可组合性质.

\section{余米田}

照例, 把箭头的方向反过来, 我们就得到一个额外的构造. 米田引理可以应用于反范畴 $\cat{C}^{op}$,
从而给出逆变函子之间的一个映射.

等价地, 我们可以通过固定 hom-函子的目标物件 (而不是源物件) 来推出余米田引理. 我们得到从
$\cat{C}$ 到 $\Set$ 的逆变 hom-函子: $\cat{C}(-, a)$. 米田引理的逆变版本确立了一一对应, 一边是
从这个函子到任意其他逆变函子 $F$ 的自然变换, 另一边是集合 $F a$ 的元素:
\[\cat{Nat}(\cat{C}(-, a), F) \cong F a\]
下面是余米田引理的 Haskell 版本:

\begin{snipv}
forall x . (x -> a) -> F x \ensuremath{\cong} F a
\end{snipv}
注意, 在一些文献里, 被称为米田引理的恰恰是这个逆变版本.

\section{挑战}

\begin{enumerate}
  \tightlist
  \item
        证明在 Haskell 中构成米田同构的两个函数 \code{phi} 与 \code{psi} 互为逆.

        \begin{snip}{haskell}
phi :: (forall x . (a -> x) -> F x) -> F a
phi alpha = alpha id

psi :: F a -> (forall x . (a -> x) -> F x)
psi fa h = fmap h fa
\end{snip}
  \item
        离散范畴是指有物件、但除了恒等态射之外再无别的态射的范畴. 从这种范畴出发的函子,
        米田引理是如何成立的?
  \item
        单位值的列表 \code{{[}(){]}} 除了长度之外不包含其他信息. 所以作为一种数据类型, 它可以
        看作整数的编码. 空列表编码零, 单例 \code{{[}(){]}} (是一个值, 不是类型) 编码一, 以此
        类推. 利用针对列表函子的米田引理, 构造出这种数据类型的另一种表示.
\end{enumerate}

\section{参考资料}

\begin{enumerate}
  \tightlist
  \item
        The Catsters 关于
        \urlref{https://www.youtube.com/watch?v=TLMxHB19khE}{米田引理}的视频.
\end{enumerate}
